\subsection{初項依存型}\label{節：初項依存型}
  \BookFigCobweb

  図では$f(x)=\displaystyle\frac{x^2-1}{2}$として,
  横軸上の$a_1=2$から曲線$y=f(x)$へ縦に進み,
  直線$y=x$へ横に進む操作を繰り返している.
  縦の移動は$a_{n+1}=f(a_n)$という計算そのものであり,
  横の移動で得られた値を次の入力として読む.
  最初の三つの入力は$a_1=2,\,a_2=\displaystyle\frac32,\,
  a_3=\displaystyle\frac58$に対応する.
  初項を変えれば出発点が変わり,同じ更新規則でも別の軌道になる.



最後に,\bookterm{recurrence}{漸化式}の形だけでなく,初項の選び方そのものが問題を左右する
例を見よう.
\[
  a_1=c,
  \quad
  a_{n+1}=\frac{a_n^2-1}{2}
\]
初項を記号 $c$ のままにすると,
\begin{align*}
  a_2&=\frac{c^2-1}{2},\\
  a_3&=\frac{c^4-2c^2-3}{8},\\
  a_4&=\frac{c^8-4c^6-2c^4+12c^2-55}{128}
\end{align*}
一般に $a_n$ は $c$ の次数 $2^{n-1}$ の多項式になる.初項を一般の
文字のまま扱おうとすると,項を一つ進めるたびに次数が倍になり,
\bookterm{trigonometric}{三角関数}型や基本対称式型のような固定された表現が見えてこない.

特に $c=2$ と固定すると,
\[
  a_1=2,
  \quad
  a_2=\frac{3}{2},
  \quad
  a_3=\frac{5}{8},
  \quad
  a_4=-\frac{39}{128}
\]
となる.しかし,この列は直前の例のように角度を2倍する形にも,
指数の差を保つ形にも,接線近似の形にもならない.

このように$c=2$では,ここまでに試した変換から直ちに簡単な表現を見出せない.ただし,数項の計算だけで,他の初項についても簡単な表現が存在しないとは結論できない.
一方で,初項をうまく選べば特別な構造が現れることもある.
例えば定\bookterm{sequence}{数列}を探すだけなら,
\[
  c=1+\sqrt{2},
  \quad
  c=1-\sqrt{2}
\]
では $a_{n+1}=a_n=c$ となる.大切なのは,ある初項を選んだときだけ
特別な恒等式や置換が働くことがあり,初項を一般化するとその偶然の
構造が失われる,という事実である.

初項によって解法の見通しが変わる例は,先に挙げた基本対称式型にもある.
もちろん,一般の初項での解法は先に示したとおりであるが,より狭義の解き方ではこのようなものがある.
\[
  a_1=-1,\quad a_{n+1}=a_n^2-2
\]
のような場合であれば,第5項までが以下のようになる.
\[
  -1,-1,-1,-1,-1
\]
\bookterm{recurrence}{漸化式}が表す\bookterm{sequence}{数列}が定数\bookterm{sequence}{数列}を成したのだ.
このような解はシンプルに\bookterm{fixed}{不動点}を求めることで作り出せる.
つまりは
\[
  \alpha=\alpha^2-2
\]
を$\alpha$について解けばよい.
この二次方程式の解は$\alpha=-1,2$である.
\,$\alpha=2$を初項に選ぶと,\bookterm{recurrence}{漸化式}から常に$a_n=2$となる.

実は,この定\bookterm{sequence}{数列}$a_n=-1$は,\,\sectionref{節：基本対称式型}で扱った$a_{n+1}=a_n^2-2$型の特別な場合と見ることができる.
\,$\alpha+\alpha^{-1}=a_1$を満たす$\alpha$が単位円上,すなわち$\alpha=e^{i\theta}$となるように($-2\le a_1\le2$のように)初項を選べば,
\[
  a_n=\alpha^{2^{n-1}}+\alpha^{-2^{n-1}}=2\cos(2^{n-1}\theta)
\]
という\bookterm{trigonometric}{三角関数}の形になる.今回の$a_1=-1$はまさにこの場合で,\,$2\cos\theta=-1$すなわち$\theta=\dfrac{2\pi}3$ととれる.このとき$2^{n-1}\theta$は$2$倍するごとに$\dfrac{2\pi}3$と$\dfrac{4\pi}3$を行き来し(どちらも余弦は$-\dfrac12$),\,$a_n=2\cos(2^{n-1}\theta)=-1$が常に成り立つ.つまり$a_1=-1$は,角度を$2$倍しても余弦の値が変わらないという特別な初項を選んだ場合に他ならない.

この例では,少なくとも本節で扱った初等的な道具からは\bookterm{general}{一般項}を得る
決まった方法が出てこない.それでも\bookterm{recurrence}{漸化式}は\bookterm{sequence}{数列}を定めており,
\[
  a_n=f^{\,\circ(n-1)}(c),
  \quad
  f(x)=\frac{x^2-1}{2}
\]
と書くことはできる.これは解答の終わりではなく,「\bookterm{general}{一般項}を求める」
という目標自体を見直す入口である.付録では,閉じた式が得られない
場合にも,軌道,\bookterm{limit}{極限},近似,\bookterm{matrix}{行列},\bookterm{generating}{母関数}など別の言葉で\bookterm{sequence}{数列}を理解する
方法を考える.
